% ------------------------------------------------------------
% A6.x Substrate Uplift Module F — Composite Substrate Algebra (v0.1)
% Author: Borealis S. Hedling
% Altitude: A6.x (Substrate Uplift Layer)
% Purpose: Define the composite substrate algebra that unifies
%          Modules A–E (RSVP, TARTAN, Unistochastic, Sheaf-
%          Cognition, Trajectory Invariants) into a single A6.x
%          altitude algebra. This is the final A6.x uplift module.
% ------------------------------------------------------------

\documentclass[11pt]{article}
\usepackage{amsmath, amssymb, amsthm}

\title{A6.x Substrate Uplift Module F: Composite Substrate Algebra}
\author{Borealis S. Hedling}
\date{October 2026}

\begin{document}
\maketitle

\section*{Overview}

This module defines the composite substrate algebra at A6.x altitude. It unifies
the symplectic RSVP substrate (Module A), the TARTAN projection algebra (Module
B), the unistochastic emergence substrate (Module C), the sheaf-cognition gluing
substrate (Module D), and the economic trajectory invariants (Module E) into a
single algebraic structure. This composite algebra is the final A6.x uplift
module and prepares the manifold for A7.1 Dual-Compatibility Transition
Operators.

% ------------------------------------------------------------
\section{Composite Substrate Space}

Define the composite substrate space:
\[
\mathcal{A}_{\mathrm{comp}} =
\mathcal{A}_{\mathrm{RSVP}}
\oplus
\mathcal{A}_{\mathrm{TARTAN}}
\oplus
\mathcal{A}_{\mathrm{uni}}
\oplus
\mathcal{A}_{\mathrm{sheaf}}
\oplus
\mathcal{A}_{\mathrm{traj}}.
\]

Each component is defined in Modules A–E.

% ------------------------------------------------------------
\section{Composite Operator}

Define the composite operator:
\[
\mathcal{O}(p) =
\mathcal{C}(p)
\oplus
\mathcal{D}(p)
\oplus
\mathcal{S}(p)
\oplus
\mathcal{X}(p)
\oplus
\mathcal{E}(p),
\]
where each term is the corresponding operator from A7.0.

\textbf{Theorem (Closure).}
\[
\mathcal{O}(p) \in \mathcal{A}_{\mathrm{comp}}.
\]

Thus the composite operator is closed under the composite substrate algebra.

% ------------------------------------------------------------
\section{Composite Lie Algebra}

Define the composite Lie bracket:
\[
[\mathcal{O}_1, \mathcal{O}_2]
=
[\mathcal{C}_1, \mathcal{C}_2]
\oplus
[\mathcal{D}_1, \mathcal{D}_2]
\oplus
[\mathcal{S}_1, \mathcal{S}_2]
\oplus
[\mathcal{X}_1, \mathcal{X}_2]
\oplus
[\mathcal{E}_1, \mathcal{E}_2].
\]

\textbf{Theorem (Lie Compatibility).}
\[
[\mathcal{O}_1, \mathcal{O}_2] \in \mathcal{A}_{\mathrm{comp}}.
\]

This enforces Fix 3 (covariant Lie bracket).

% ------------------------------------------------------------
\section{Composite Spectral Algebra}

Define:
\[
\mathrm{Spec}(\mathcal{O}(p)) =
\mathrm{Spec}(\mathcal{C}(p))
\cup
\mathrm{Spec}(\mathcal{S}(p)).
\]

\textbf{Theorem (Spectral Continuity).}
If each component spectrum is continuous, then:
\[
\mathrm{Spec}(\mathcal{O}(p)) \text{ continuous in } p.
\]

This enforces Fix 5 and Fix 6.

% ------------------------------------------------------------
\section{Composite Modal Algebra}

Define:
\[
\mathcal{M}_{\mathrm{comp}}(p) =
\mu(\Pi_{\mathrm{T}}(p))
\cup
\mu(\mathcal{A}_{\mathrm{sheaf}}(p)).
\]

\textbf{Theorem (Modal Lift).}
If modal truth lifts through each component, then:
\[
\mathcal{M}_{\mathrm{comp}}(p) \leadsto \mathcal{M}_{\mathrm{comp}}(q).
\]

This enforces Fix 10.

% ------------------------------------------------------------
\section{Composite Provenance Algebra}

Define:
\[
\mathcal{P}_{\mathrm{comp}}(p) =
I_{\mathrm{prov}}(\Pi_{\mathrm{T}}(p))
\cup
I_{\mathrm{prov}}(\mathcal{A}_{\mathrm{sheaf}}(p)).
\]

\textbf{Theorem (Provenance Lift).}
If provenance lifts through each component, then:
\[
\mathcal{P}_{\mathrm{comp}}(M) \Rightarrow
\mathcal{P}_{\mathrm{comp}}(\mathcal{A}(M,S)).
\]

This enforces Fix 11.

% ------------------------------------------------------------
\section{Composite Semantic Algebra}

Define:
\[
\mathcal{V}_{\mathrm{comp}}(p) =
I_{\mathrm{sem}}(\Pi_{\mathrm{T}}(p))
\cup
I_{\mathrm{sem}}(\mathcal{A}_{\mathrm{sheaf}}(p)).
\]

\textbf{Theorem (Semantic Lift).}
If semantic truth lifts through each component, then:
\[
\mathcal{V}_{\mathrm{comp}}(p) \leadsto \mathcal{V}_{\mathrm{comp}}(q).
\]

This enforces Fix 12.

% ------------------------------------------------------------
\section{Collapse-Surface Compatibility}

Define:
\[
\Sigma_{\mathrm{comp}} =
\Sigma_{\mathrm{RSVP}}
\cup
\Sigma_{\mathrm{TARTAN}}
\cup
\Sigma_{\mathrm{uni}}
\cup
\Sigma_{\mathrm{sheaf}}
\cup
\Sigma_{\mathrm{traj}}.
\]

\textbf{Theorem (Local Detectability).}
If collapse is locally detectable in each component, then:
\[
\Sigma_{\mathrm{comp}} \text{ locally detectable}.
\]

This enforces Fix 8.

% ------------------------------------------------------------
\section{Conclusion}

The composite substrate algebra defined here unifies all A6.x substrate modules
into a single algebraic structure. This is the final A6.x uplift module and
prepares the manifold for A7.1 Dual-Compatibility Transition Operators.

\end{document}
